\chapter{Usability Test Instrument}
\label{app:sus-questionnaire}

This appendix reproduces the System Usability Scale (SUS), the
standardized 10-item questionnaire used in \cref{sec:usability-eval}, in
its original form~\cite{brooke1996sus}. Each item is rated on a 5-point
Likert scale from \emph{Strongly disagree} (1) to \emph{Strongly agree}
(5). The overall score is computed by summing, for odd-numbered items,
(rating $-$ 1), and for even-numbered items, (5 $-$ rating), then
multiplying the total by 2.5 to obtain a score in $[0,100]$.

\begin{enumerate}
  \item I think that I would like to use this app frequently.
  \item I found the app unnecessarily complex.
  \item I thought the app was easy to use.
  \item I think that I would need the support of a technical person to be
  able to use this app.
  \item I found the various functions in this app were well integrated.
  \item I thought there was too much inconsistency in this app.
  \item I would imagine that most people would learn to use this app very
  quickly.
  \item I found the app very cumbersome/awkward to use.
  \item I felt very confident using the app.
  \item I needed to learn a lot of things before I could get going with
  this app.
\end{enumerate}

\section*{Post-task open question}

In addition to the SUS instrument, each participant was asked the
following scenario-specific question immediately after completing the
report-creation task, directly targeting NFR-5 (sync-status visibility,
\cref{sec:nonfunctional-req}):

\begin{quote}
\emph{"At any point while completing this task, were you unsure whether
your report had been saved, or whether it had been sent anywhere?"}
\end{quote}

Responses should be recorded verbatim and coded as \emph{clear} (the
participant could state, unprompted, that the report was saved locally
and would sync later) or \emph{unclear} (the participant expressed doubt
or could not explain the report's state).
